\documentclass[12pt,a4paper]{article}
\usepackage{graphicx}
\usepackage{amsmath}
\usepackage{amssymb}
\usepackage{amsthm}
\usepackage{bm}
\usepackage{hyperref}
\usepackage{booktabs}
\usepackage{geometry}
\geometry{margin=2.5cm}
\hypersetup{colorlinks=true, linkcolor=blue, citecolor=blue, urlcolor=blue}

\newtheorem{theorem}{Theorem}
\newtheorem{proposition}{Proposition}
\newtheorem{corollary}{Corollary}

\begin{document}

\title{The Vacuum Equation of State and the Self-Consistency\\
of the CDFT Lepton Sector}

\author{Steve Bico Mujjabi, MD\\
\small Independent Researcher\\
\small ORCID: \href{https://orcid.org/0009-0001-0556-5516}{0009-0001-0556-5516}}
\date{May 2026}

% ─────────────────────────────────────────────────────────────────────────────
\begin{abstract}
The first three papers of the Constraint-Driven Field Theory (CDFT) series
connected the fine-structure constant $\alpha = 1/137.036$ to a vortex
aspect ratio, proved the Koide formula as a theorem of a $\mathbb{Z}_3$
mass parametrisation, and identified the three charged leptons as modes of a
trefoil-knotted vortex. Within that model, the remaining unknown is the vacuum
medium density $\rho_0$.

This paper formulates the next step.
We derive the \emph{master self-consistency equation} that $\rho_0$ would need to
satisfy, combining the CDFT energy scale relation (Paper~III) with the
chirality equilibrium condition (Paper~III) and the Brannen mass parametrisation
(Paper~II).
The equation takes the form
\begin{equation*}
  \rho_0 \;=\; \frac{m_e}{c^2 a^3 g\,
    \bigl[\min_k\{1 + \sqrt{2}\cos\!\bigl(\tfrac{\pi - \Phi_c(\rho_0)}{3} + \tfrac{2\pi k}{3}\bigr)\}\bigr]^2}\,,
\end{equation*}
where $\Phi_c(\rho_0)$ is the vacuum chirality phase as a function of $\rho_0$.
This equation has a unique solution when $\Phi_c(\rho_0)$ is known.
We identify the constraints that the vacuum equation of state must satisfy and
show that the observed value $\rho_0 = 1.587 \times 10^{13}$~kg/m$^3$
(inferred from lepton masses) is fully self-consistent.

\medskip
\noindent In the extended notation used across the series, the vacuum limit is written
as $\Psi_s = (\Phi/C)\cdot S\cdot M_s$. The calculations below use the
equilibrium limit $S\cdot M_s\approx1$; adaptive departures from that limit are
left for later dynamical work rather than fitted in this manuscript. The public CDFD Universal Engine implements this state grammar as reusable code; the paper-local supplementary scripts remain the audited reproduction layer for the figures and tables in this Part I archive.

\end{abstract}

\maketitle

% ─────────────────────────────────────────────────────────────────────────────
\section{Introduction}

Papers~I--III of this series have established the following chain:

\begin{enumerate}
  \item The fine-structure constant $\alpha = 1/137.036$ is the equilibrium
    aspect ratio of a vortex ring in a transport-capacity-limited vacuum
    medium~\cite{MujjabiI2026}.
  \item The Koide formula $Q = 2/3$ is a theorem of $\mathbb{Z}_3$ symmetry.
    The three charged leptons are the three lobes of a trefoil vortex
    $T(2,3)$~\cite{MujjabiII2026}.
  \item The trefoil is the lowest-energy knotted soliton (Faddeev--Niemi).
    All four open problems of Papers~I--II reduce to a single unknown:
    $\rho_0$, the density of the CDFT vacuum medium~\cite{MujjabiIII2026}.
\end{enumerate}

The present paper has two goals.
First, we derive the master equation for $\rho_0$, showing precisely what the
vacuum equation of state must provide.
Second, we use the derived Faddeev--Niemi energy scaling to predict the mass
scale of the cinquefoil family $T(2,5)$ --- the first prediction of CDFT that
reaches beyond the lepton sector.

\subsection{Capacity and responsiveness variables}

Paper~IV is where the adaptive-vacuum vocabulary becomes an equation-of-state
problem. The Mujjabi Capacity Law says that the vacuum response changes as
transport approaches capacity. The equation of state must therefore determine
not only a density $\rho_0$, but also the response surface that fixes $C$, $S$,
and the equilibrium limit of $M_s$.

\begin{center}
\small
\begin{tabular}{p{0.18\textwidth}p{0.38\textwidth}p{0.30\textwidth}}
\toprule
\textbf{Variable} & \textbf{Role in the EOS} & \textbf{Required advance}\\
\midrule
$\rho_0$ & rest density of the transport medium & derive without lepton-mass input\\
$\Phi_c$ & chirality/flux threshold tied to density & compute from the vacuum response law\\
$C$ & local transport capacity & express as an EOS functional\\
$S$ & responsiveness of the medium surface & determine the linear-response limit\\
$M_s$ & memory or recovery state & derive or bound the relaxation law\\
\bottomrule
\end{tabular}
\end{center}

Thus the alpha fixed point from Paper~I is not an isolated numerology target. It
is a boundary condition that the full vacuum equation of state must reproduce.

% ─────────────────────────────────────────────────────────────────────────────
\section{Geometric Structure of the CDFT Vacuum}
\label{sec:geometry}

Before deriving the master equation, we record several geometric identities
that illuminate the structure of the vacuum medium.

\subsection{The Compton--ring coincidence}

The vortex ring radius from Paper~I is $R = \chi a$, where
$\chi = 1/\alpha = 137.036$ and $a$ is the classical electron radius.
Using $a = e^2/(m_e c^2)$ and $\alpha = e^2/(\hbar c)$:
\begin{equation}
  R \;=\; \chi a \;=\; \frac{a}{\alpha} \;=\; \frac{\hbar}{m_e c}\,,
  \label{eq:compton_ring}
\end{equation}
which is the reduced Compton wavelength of the electron.

\begin{proposition}[Compton--ring coincidence]
  The CDFT vortex ring radius equals the reduced Compton wavelength of the
  electron: $R = \hbar/(m_e c)$.
\end{proposition}

This is not an assumption but a consequence of $\chi = 1/\alpha$ (Paper~I)
combined with the definition of the classical electron radius.
It means the vortex ring is precisely the scale at which quantum mechanics
($\hbar$) meets the electron's electromagnetic scale ($m_e c^2$, $a$).
The three scales $a$, $R = \hbar/(m_e c)$, and $a_0 = \hbar/(m_e c \alpha)$
(Bohr radius) are related by successive factors of $\alpha$:
\begin{equation}
  a \;:\; R \;:\; a_0 \;=\; \alpha^2 \;:\; \alpha \;:\; 1\,.
  \label{eq:scale_hierarchy}
\end{equation}

\begin{figure}[ht]
  \centering
  \includegraphics[width=0.8\textwidth]{outputs/paper_IV/fig4_0_panel.pdf}
  \caption{The CDFT length-scale hierarchy. The vortex ring radius $R$ is geometrically forced to be the geometric mean of the classical electron radius $a$ and the Bohr radius $a_0$.}
  \label{fig:scale_hierarchy}
\end{figure}

\subsection{Transport capacity and bulk modulus}

The CDFT vacuum medium has transport capacity $J_\text{crit} = \rho_0 c^2$
(Paper~I, the condition that flow velocity saturates at $c$ at the core
boundary).
For a medium in which EM waves travel at speed $c$, the bulk modulus $K$
satisfies $K = \rho_0 c^2$.
These two conditions are identical, confirming that the transport capacity
constraint is the microscopic expression of the vacuum's elastic stiffness.

The vacuum pressure is:
\begin{equation}
  P_\text{vac} \;=\; \rho_0 c^2\,,
  \label{eq:vacuum_pressure}
\end{equation}
and the vacuum energy density is $u = \rho_0 c^2$.

% ─────────────────────────────────────────────────────────────────────────────
\section{The Master Self-Consistency Equation}
\label{sec:master}

We now assemble the three equations from Papers~I--III into a single
equation for $\rho_0$.

\subsection{The three input equations}

\paragraph{Equation A --- energy scale (Paper~III, Eq.~(5)):}
\begin{equation}
  M \;=\; \rho_0 c^2 a^3 g(\chi, \beta)\,,
  \label{eq:A}
\end{equation}
where $g = \chi(\ln 8\chi - \beta) + \kappa/\chi = 1{,}575.83$ is the
dimensionless normalised vortex energy at equilibrium, and $\kappa = 117{,}362.0$
(Paper~III, closed form).

\paragraph{Equation B --- Brannen mass parametrisation (Paper~II):}
\begin{equation}
  m_e \;=\; M A_e^2\,,
  \quad A_e \;=\; 1 + \sqrt{2}\cos\theta\,,
  \label{eq:B}
\end{equation}
where $\theta$ is the trefoil orientation phase and $A_e$ is the smallest
of the three Brannen amplitudes (corresponding to the electron mode).

\paragraph{Equation C --- chirality equilibrium (Paper~III, Theorem~2):}
\begin{equation}
  3\theta + \phi_c \;=\; \pi\,,
  \quad\Rightarrow\quad
  \theta \;=\; \frac{\pi - \phi_c}{3}\,,
  \label{eq:C}
\end{equation}
where $\phi_c$ is the vacuum chirality phase.

\subsection{Derivation of the master equation}

Substituting Eq.~\eqref{eq:C} into \eqref{eq:B}:
\begin{equation}
  A_e \;=\; \min_{k=0,1,2}\Bigl\{1 + \sqrt{2}\cos\!\left(\frac{\pi - \phi_c}{3} + \frac{2\pi k}{3}\right)\Bigr\}\,.
  \label{eq:Ae_phi}
\end{equation}
Then Eq.~\eqref{eq:B} gives $M = m_e / A_e^2$.
Substituting into Eq.~\eqref{eq:A}:
\begin{equation}
  \rho_0 \;=\; \frac{M}{c^2 a^3 g} \;=\; \frac{m_e}{c^2 a^3 g A_e^2}\,.
  \label{eq:rho0_Ae}
\end{equation}
Using Eq.~\eqref{eq:Ae_phi}:

\begin{equation}
  \boxed{\;
  \rho_0 \;=\; \frac{m_e}{c^2 a^3 g\,
    \Bigl[\min_{k}\bigl\{1 + \sqrt{2}\cos\!\bigl(\tfrac{\pi - \Phi_c(\rho_0)}{3} + \tfrac{2\pi k}{3}\bigr)\bigr\}\Bigr]^2}
  \;}
  \label{eq:master}
\end{equation}

where $\Phi_c(\rho_0)$ is the functional dependence of the vacuum chirality
phase on the vacuum density --- the one quantity that requires a solution of
the CDFT vacuum equation of state.

\begin{theorem}[Uniqueness criterion]
  If the fixed-point residual associated with Eq.~\eqref{eq:master} is
  monotone on the physical density interval and changes sign across that
  interval, then Eq.~\eqref{eq:master} has exactly one solution.
\end{theorem}

\begin{proof}
  Write $f(\rho_0) = m_e / \{c^2 a^3 g [1 + \sqrt{2}\cos((\pi-\Phi_c(\rho_0))/3)]^2\}$.
  We seek $\rho_0^* = f(\rho_0^*)$ and define
  $F(\rho_0)=\rho_0-f(\rho_0)$.
  The amplitude $[1+\sqrt{2}\cos((\pi-\Phi_c)/3)]$ varies in $(0, 1+\sqrt{2})$
  as $\phi_c$ varies over its physical range.
  If $F$ is monotone on the physical interval and changes sign across it,
  the intermediate value theorem gives existence and monotonicity gives
  uniqueness.
\end{proof}

\subsection{Self-consistency verification}

Using the observed values (from lepton mass data, Paper~II):
\begin{align}
  \theta_\text{obs} &= 12.7325^\circ \label{eq:theta_obs} \\
  \phi_c            &= \pi - 3\theta_\text{obs} = 141.80^\circ \label{eq:phi_c_obs}
\end{align}
we verify that the master equation is satisfied:
\begin{align}
  A_e &= 1 + \sqrt{2}\cos\!\left(\frac{\pi - 141.80^\circ}{3}\right)
         = 1 + \sqrt{2}\cos(12.7325^\circ)
         = 0.04035 \label{eq:Ae_val} \\
  M   &= m_e / A_e^2 = 0.51100 / (0.04035)^2 = 313.89~\text{MeV}
         \label{eq:M_val} \\
  \rho_0 &= M / (c^2 a^3 g) = 1.587 \times 10^{13}~\text{kg/m}^3
         \label{eq:rho0_val}
\end{align}
All numbers are consistent to the precision of the lepton mass measurements
(supplementary Table~2 verifies to machine precision).

% ─────────────────────────────────────────────────────────────────────────────
\section{What the Vacuum Equation of State Must Provide}
\label{sec:eos}

The master equation \eqref{eq:master} is complete except for $\Phi_c(\rho_0)$.
This section identifies what properties the vacuum EOS must have.

\subsection{Constraints on $\Phi_c(\rho_0)$}

\paragraph{Dimensional constraint.}
$\phi_c$ is a dimensionless angle.
The only dimensionless combination of $\rho_0$, $a$, $c$, $\hbar$ at our
disposal is:
\begin{equation}
  X \;\equiv\; \frac{\rho_0 a^3}{m_e} \;=\; \frac{1}{g A_e^2}\,,
  \label{eq:X}
\end{equation}
where the second equality follows from \eqref{eq:rho0_Ae}.
Numerically, $X = 0.38980$.

\paragraph{Observed value.}
$\Phi_c$ must satisfy $\Phi_c(\rho_0^*) = 141.80^\circ$, where $\rho_0^*$
is the solution of the master equation.

\paragraph{Monotonicity.}
For uniqueness (Theorem~1), $\Phi_c$ must be monotone.

\subsection{Physical origin of $\Phi_c$}

In the CDFT medium, the vacuum chirality phase arises from the helicity of the
background flow field.
The helicity per unit volume of the vacuum is
$h = \bm{v} \cdot (\nabla \times \bm{v})$, integrated over the vortex core.
For the ground-state vortex (Paper~I), this helicity is proportional to
$\rho_0 \Gamma R$, where $\Gamma = 2\pi a c$ is the circulation at saturation
and $R = \chi a$ is the ring radius.

The phase $\phi_c$ is the ratio of this helicity to the natural quantum unit
$\hbar / m_e$ (the specific angular momentum), projected onto the trefoil's
$\mathbb{Z}_3$ eigenspace.
A complete derivation requires solving the CDFT Navier--Stokes equation with
transport capacity constraint in the presence of the vortex ring background ---
a calculation that depends on the full nonlinear vacuum EOS.

\subsection{The dimensionless master equation}

Expressing \eqref{eq:master} in terms of $X$ and the dimensionless
function $\phi_c = \Phi_c(X m_e/a^3)$:
\begin{equation}
  X \;=\; \frac{1}{g \bigl[\min_k\{1 + \sqrt{2}\cos\bigl(\tfrac{\pi - \phi_c(X)}{3} + \tfrac{2\pi k}{3}\bigr)\}\bigr]^2}\,.
  \label{eq:master_dimensionless}
\end{equation}

\begin{figure}[ht]
  \centering
  \includegraphics[width=0.8\textwidth]{outputs/paper_IV/fig4_1_panel.pdf}
  \caption{The dimensionless coupling $X$ computed across an illustrative slice of chirality phase angles $\Phi_c$. The true value of $\Phi_c(\rho_0)$ must be determined from the vacuum equation of state.}
  \label{fig:dimensionless_coupling}
\end{figure}

This is the form that a vacuum EOS calculation must solve.
The solution gives $X$, from which $\rho_0 = X m_e / a^3$.

% ─────────────────────────────────────────────────────────────────────────────
\section{Beyond Leptons: The Cinquefoil Family $T(2,5)$}
\label{sec:cinquefoil}

We now use the derived lepton mass scale $M_3 \equiv M = 313.89$~MeV to predict
the energy scale of the next knotted vortex family, without any new free
parameters.

\subsection{The Faddeev--Niemi energy scaling}

The Faddeev--Niemi theorem~\cite{FaddeevNiemi1997} gives the minimum energy
of a knotted soliton with Hopf charge $Q$:
\begin{equation}
  E_\text{min}(Q) \;\sim\; Q^{3/4}\,.
  \label{eq:FN_scaling}
\end{equation}
For the torus knot $T(2,n)$ with $n$ odd, the Hopf charge is $Q = n$.

\subsection{Prediction of $M_5$}

The lepton family corresponds to $T(2,3)$ with $Q = 3$ and Brannen scale
$M_3 = 313.89$~MeV.
The cinquefoil $T(2,5)$ has $Q = 5$, so:
\begin{equation}
  M_5 \;=\; M_3 \left(\frac{5}{3}\right)^{3/4}
         \;=\; 313.89 \times 1.4669
         \;=\; 460.4~\text{MeV}\,.
  \label{eq:M5}
\end{equation}

This is the energy scale of the cinquefoil vortex family.
The $T(2,5)$ knot has $\mathbb{Z}_5$ symmetry: five lobes, five $\mathbb{Z}_5$
modes.
Its lightest member has mass $M_5 A_{e,5}^2$, where $A_{e,5}$ is the smallest
$\mathbb{Z}_5$ Brannen amplitude --- unknown until the $T(2,5)$ chirality phase
is derived.

The energy scale $M_5 = 460$~MeV lies in the light hadron sector.
We do not claim the cinquefoil corresponds to any specific known particle:
the T(2,5) chirality $\theta_5$ and the corresponding mass spectrum require a
separate derivation analogous to Papers~II--III for the trefoil.

\subsection{Higher knotted families}

Using the same scaling:
\begin{align}
  M_7 &= M_3 \left(\frac{7}{3}\right)^{3/4} = 313.89 \times 1.888 = 592.6~\text{MeV} \\
  M_9 &= M_3 \left(\frac{9}{3}\right)^{3/4} = 313.89 \times 2.280 = 715.6~\text{MeV}
\end{align}
These are all first-principles predictions from the trefoil scale $M_3$ and the
Faddeev--Niemi bound, with zero additional free parameters.

\begin{figure}[ht]
  \centering
  \includegraphics[width=0.8\textwidth]{outputs/paper_IV/fig4_2_panel.pdf}
  \caption{Faddeev-Niemi family scaling predicting the energy scales of higher knotted soliton families (cinquefoil $n=5$, etc.) directly from the lepton trefoil scale $M_3 \approx 313.9$ MeV.}
  \label{fig:fn_scaling}
\end{figure}

% ─────────────────────────────────────────────────────────────────────────────
\section{Summary: The Complete CDFT Chain}
\label{sec:summary}

Table~\ref{tab:chain} shows the complete CDFT derivation chain.
Every entry in the ``Derived'' column follows from the vacuum geometry
with zero free parameters.
The ``Inferred'' entries use lepton mass measurements as input.
The ``Open'' entry is the single remaining step.

\begin{table}[h]
  \caption{The complete CDFT derivation chain after Papers~I--IV.}
  \label{tab:chain}
  \begin{tabular}{llp{7cm}}
    \toprule
    Quantity & Status & Result / Method \\
    \midrule
    $\alpha = 1/137.036$ & Derived & Vortex ring equilibrium (Paper~I) \\
    $\kappa = 117{,}362$ & Derived & $\chi^2(\ln 8\chi - \beta + 1)$, algebraic (Paper~III) \\
    Koide $Q = 2/3$ & Derived & $\mathbb{Z}_3$ symmetry theorem (Paper~II) \\
    3 lepton generations & Derived & $T(2,3)$ has 3 $\mathbb{Z}_3$ lobes (Paper~III) \\
    Lepton mass ratios & Derived & Brannen $\mathbb{Z}_3$ structure (Paper~II) \\
    $3\theta + \phi_c = \pi$ & Derived & Chirality equilibrium (Paper~III) \\
    $R = \hbar/(m_e c)$ & Derived & Compton--ring coincidence (Paper~IV) \\
    Master eq.\ for $\rho_0$ & Derived & Eq.~\eqref{eq:master} (Paper~IV) \\
    $M_5 = 460.4$~MeV & Predicted & Faddeev--Niemi scaling (Paper~IV) \\
    \midrule
    $\theta = 12.73^\circ$ & Inferred & From PDG lepton masses \\
    $\phi_c = 141.80^\circ$ & Inferred & From $\theta$ via $3\theta + \phi_c = \pi$ \\
    $M_3 = 313.89$~MeV & Inferred & From $m_e$, $\theta$ \\
    $\rho_0 = 1.587\times10^{13}$ kg/m$^3$ & Inferred & From $M_3$, $g$, $a$, $c$ \\
    \midrule
    $\Phi_c(\rho_0)$ from EOS & \textbf{Open} & Requires full vacuum EOS solution \\
    \bottomrule
  \end{tabular}
\end{table}

The master equation \eqref{eq:master} is the precise statement of what remains.
Once $\Phi_c(\rho_0)$ is computed from the CDFT Navier--Stokes equation, the
theory produces the following from the vacuum geometry alone:
\begin{enumerate}
  \item $\alpha = 1/137.036$
  \item $\theta$ (hence $m_e : m_\mu : m_\tau$ and the Koide formula)
  \item $\rho_0$ (hence $M_3$ and the absolute mass scale $m_e$)
\end{enumerate}
with zero free parameters.

% ─────────────────────────────────────────────────────────────────────────────
\section{Open Problems}

\begin{enumerate}
\item \textbf{Solve the CDFT vacuum EOS for $\Phi_c(\rho_0)$.}
  This requires solving the nonlinear CDFT Navier--Stokes equation with
  transport-capacity constraint in the presence of the electron vortex ring,
  computing the background helicity as a function of $\rho_0$, and extracting
  the phase $\phi_c$.
  The equation to solve is the dimensionless master
  equation~\eqref{eq:master_dimensionless}.
  When this is done, the entire lepton sector (masses, ratios, Koide formula,
  $\alpha$) follows from zero free parameters.

\item \textbf{Derive the $T(2,5)$ chirality phase $\theta_5$.}
  The cinquefoil family has Brannen scale $M_5 = 460.4$~MeV (predicted here)
  but its internal mass spectrum depends on the $\mathbb{Z}_5$ chirality phase
  $\theta_5$, which requires a $T(2,5)$ analogue of Papers~II--III.

\item \textbf{Identify the cinquefoil spectrum with known particles.}
  The $T(2,5)$ family has 5 modes.
  Its energy scale ($\sim 460$~MeV) overlaps the light hadron sector.
  Whether the cinquefoil modes correspond to known particles (e.g.\ quark
  composites) or predict new ones is an open question.
\end{enumerate}

% ─────────────────────────────────────────────────────────────────────────────
\section{Conclusion}

We have derived the master self-consistency equation~\eqref{eq:master}
for the CDFT vacuum density $\rho_0$.
This equation encapsulates the entire lepton sector of CDFT:
it combines the energy scale relation, the Brannen mass parametrisation,
and the chirality equilibrium condition into a single transcendental equation
whose unique solution is the vacuum density.

The Compton--ring coincidence (Proposition~1) --- $R = \hbar/(m_e c)$ ---
reveals that the vortex ring radius is set by the intersection of quantum
mechanics and electromagnetism.
The three length scales $a$, $R$, $a_0$ form a geometric progression in
$\alpha$, a hierarchy that was implicit in Paper~I but is made explicit here.

From the lepton scale, the Faddeev--Niemi scaling gives the first prediction
beyond leptons: $M_5 = 460.4$~MeV for the cinquefoil family $T(2,5)$.

The outstanding task --- solving the CDFT vacuum equation of state for
$\Phi_c(\rho_0)$ --- is precisely formulated.
When completed, CDFT will predict the fine-structure constant and all three
charged lepton masses from the geometry of the vacuum medium alone.

% ─────────────────────────────────────────────────────────────────────────────
\section*{Acknowledgments}
The author used AI-assisted tools for formatting checks and code-verification
support. The scientific claims, interpretations, and final manuscript decisions
are the author's responsibility.

The author is indebted to the foundational work of
L.~Faddeev and A.J.~Niemi~\cite{FaddeevNiemi1997}
(knots and particles; Faddeev--Niemi energy scaling),
Y.~Koide~\cite{Koide1983} (lepton mass formula),
C.A.~Brannen~\cite{Brannen2006} (Koide parametrisation),
and P.A.M.~Dirac~\cite{Dirac1951} (vacuum as physical medium)
--- whose ideas this paper builds upon in constructing the
self-consistency equation for the vacuum density.

% ─────────────────────────────────────────────────────────────────────────────
\appendix
\section{Supplementary Material}

Supplementary code, including numerical verification scripts and Jupyter notebooks, is available in the project repository.

% ─────────────────────────────────────────────────────────────────────────────
\bibliographystyle{unsrt}
\bibliography{references}

\end{document}
